\documentclass{article}
\usepackage[UTF8]{ctex}
\usepackage[tc]{titlepic}
\usepackage{amsmath, amsthm, amssymb, graphicx}
\usepackage{hyperref}
\usepackage{listings}
\usepackage{titlesec}
\usepackage{cite}
\usepackage{fancyhdr}
\usepackage{booktabs}
\usepackage{graphicx}
\usepackage{geometry}
\usepackage{caption}
\usepackage{subfigure}
\usepackage[section]{placeins}
\geometry{a4paper,scale=0.8}
\pagestyle{fancy}

\lhead{第 3 次作业\\\today}
\chead{中国科学技术大学\\数学建模课程}

\rhead{Assignment 3\\ {\CTEXoptions[today=old]\today}}
\newcommand{\upcite}[1]{\textsuperscript{\cite{#1}}}

\titleformat*{\section}{\bfseries\Large}
\titleformat*{\subsection}{\bfseries\large}

\title{\bfseries 数据拟合模型-使用神经网络进行昆虫分类}
\author{Xiaoma \quad   \quad  }

\begin{document}
\maketitle
\begin{abstract}
    神经网络作为一种已被广泛使用的模型，已经被应用于各个领域，本文将使用
    全连接神经网络探索昆虫的分类问题，并对分类结果进行讨论。
    
    
\end{abstract}

% \setcounter{secnumdepth}{1}
 \setcounter{section}{1}
\section*{\centerline{一、前言}}
人工神经网络（Artificial Neural Networks，简写为ANNs）
也简称为神经网络（NNs）或称作连接模型（Connection Model），
它是一种模仿动物神经网络行为特征，进行分布式并行信息处理的算法数学模型。
这种网络依靠系统的复杂程度，通过调整内部大量节点之间相互连接的关系，
从而达到处理信息的目的。

神经网络的优势：
\begin{itemize}
    \item 解决线性模型难以处理的线性不可分问题
    \item 线性模型泛化能力较弱
    \item 等等
\end{itemize}

全连接神经网络的结构图为
\begin{figure*}[htbp]
    \centering
	\includegraphics[width=0.8\textwidth]{./img/1.png}
\end{figure*}


 \setcounter{section}{2}
\section*{\centerline{二、相关工作}}
\begin{enumerate}
    \item 基于全连接神经网络实现昆虫分类问题
    \item 比较不同网络结构对分类器性能的影响
    \item 比较不同优化器，激活函数对分类器性能的影响
\end{enumerate}

 \setcounter{section}{3}
\section*{\centerline{三、问题分析}}
    \subsection{分析}
根据数据的聚类分布，我们发现昆虫的分类边界是非线性的，故使用线性模型很难得到
较好的结果，故使用全连接神经网络模型。
\begin{figure*}[htbp]
    \centering
	\includegraphics[width=0.6\textwidth]{./img/2.png}
\end{figure*}
\setcounter{section}{4}
\section*{\centerline{四、建模的假设}}
    \subsection{假设1}
    假设数据集中没有确实值
    \subsection{假设2}
    假设每个数据对应的标签一定正确



 \setcounter{section}{5}
 \section*{\centerline{五、符号说明}}
 \begin{table}[htbp]
    
    \centering%把表居中
    \begin{tabular}{cc}%内容全部居中
    \toprule%第一道横线
    符号&说明 \\
    \midrule%第二道横线 
    \textbf{Model} & 全连接神经网络分类器\\
    \textbf{Net} & 全连接神经网络模型\\
    \textbf{Val\_ACC} & 分类器的验证集精度\\
    \textbf{Train\_ACC} & 分类器的训练集精度\\
    \textbf{Val\_ACC} & 分类器的验证集损失\\
    \textbf{Train\_ACC} & 分类器的训练集损失\\

    \bottomrule%第三道横线
    \end{tabular}
\end{table}

 \setcounter{section}{6}
\section*{\centerline{六、数学模型建立}}
\subsection{神经网络的结构}
首先搭建一个神经网络模型，其结构为（初始结构不代表最终网络结构）
\begin{figure*}[htbp]
    \centering
	\includegraphics[width=0.8\textwidth]{./img/3.png}
\end{figure*}
\clearpage
\subsection{模型调节-网络层数与神经元个数}
已知数据集样本个数与特征数目较少，故首先假设该分类器不需要过大的网络深度与神经元个数，分别
尝试不同的网络层数与神经元个数，结果为
\begin{table}[htbp]
    
    \centering%把表居中
    \begin{tabular}{cccc}%内容全部居中
    \toprule%第一道横线
    网络层数&每层最大神经元数目&\textbf{Train\_ACC}&\textbf{Val\_ACC} \\
    \midrule%第二道横线 
    \textbf{3} & \textbf{100} &\textbf{0.935} &\textbf{0.955}\\
    \textbf{5} & \textbf{100} &\textbf{0.947} &\textbf{0.951}\\
    \textbf{10}  & \textbf{100} &\textbf{0.952} &\textbf{0.976}\\
    \textbf{3}  & \textbf{10} &\textbf{0.926} &\textbf{0.933}\\
    \textbf{5} & \textbf{10} &\textbf{0.910} &\textbf{0.922}\\
    \textbf{10}  & \textbf{10} &\textbf{0.916} &\textbf{0.948}\\

    \bottomrule%第三道横线
    \end{tabular}
\end{table}

\subsection{模型调节-优化器与激活函数}
\subsubsection{优化器选择}
在上述参数调节过程中使用的优化器为\textbf{optim.Adam}，将其训练过程的loss曲线与\textbf{optim.SGD}进
行比较
\begin{figure}[h!]
    \centering
        \subfigure[$LOSS-Adam$]{\includegraphics[width=0.4\hsize]{./img/5.png}}  
        \subfigure[$ACC-SGD$]{\includegraphics[width=0.4\hsize]{./img/7.png}} 
        
\end{figure}


可以发现使用\textbf{optim.Adam}的分类器在训练过程中收敛速度和收敛值均优于使用\textbf{optim.SGD}。
 
\subsubsection{激活函数选择}
分别使用\textbf{ReLU, Softmax, Tanh}进行分类器训练，并比较其精度的变化情况。
\begin{figure}[h!]
    \centering
        \subfigure[$ACC-ReLU$]{\includegraphics[width=0.4\hsize]{./img/4.png}}  
        \subfigure[$ACC-Softmax$]{\includegraphics[width=0.4\hsize]{./img/8.png}} 
        \subfigure[$ACC-Tanh$]{\includegraphics[width=0.4\hsize]{./img/9.png}} 
\end{figure}

发现使用\textbf{ReLU}作为激活函数时，分类器表现最优。
\setcounter{section}{7}
 \section*{\centerline{七、结果（与对比）}}
根据上述参数调节，我们得到的最优网络结构如下，
\begin{figure*}[htbp]
    \centering
	\includegraphics[width=0.8\textwidth]{./img/6.png}
\end{figure*}
使用其训练得到分类器并保存模型。
\begin{figure}[h!]
    \centering
        \subfigure[$LOSS$]{\includegraphics[width=0.4\hsize]{./img/4.png}}  
        \subfigure[$ACC$]{\includegraphics[width=0.4\hsize]{./img/5.png}} 
        
\end{figure}

 \setcounter{section}{8}
\section*{\centerline{八、结论}}
    \begin{enumerate}
        \item 若激活函数与优化器合适，增大网络深度与宽度可以提升分类器性能
        \item 若数据集较小且规整，网络结构的调整对分类器的性能影响较小
        \item 对于本次实验数据集，Adam优化比SGD优化更好
        \item 对于本次实验数据集，激活函数选择ReLU性能更好
    \end{enumerate}
 \setcounter{section}{9}
\section*{\centerline{九、问题}}
    参数调节过程为分步调参，二维交叉调参。
    
    由于时间原因，并没有尝试其他神经网络框架进行分析。




    

\bibliographystyle{plain}
\bibliography{refer}
\end{document}